\subsection{Encouraging the Development of Altruistic, Antifascist, and Ethical AI Systems}
\label{sec:9.1.4}
\runin{Policies to Counteract Authoritarian Uses of AI}
A government cannot restrict what it cannot see, and measurement belongs upstream of the whole policy toolkit for that reason. The Carnegie Endowment's AI Global Surveillance Index assessed 176 countries and found AI-enabled surveillance technology already in active use in at least 75 of them \autocite{feldstein2019globalexpansion}. That number is the evidence base for a countermeasure and is not itself one. It also demonstrates what a measurement instrument has to do to be worth the effort: count deployments instead of commitments, name the countries, and publish on a schedule that makes a change from one year to the next visible. An index reporting only that surveillance AI exists somewhere in the world would give a policymaker nothing to act on.

\runin{AI Systems That Support Democratic Values}
Section~\ref{sec:5.4.2} works through what this looks like. One gap those examples leave open is election-integrity monitoring specifically: a system built to flag anomalies in real-time reporting, or to detect coordinated interference during a vote. What would make such a system trustworthy is the list section~\ref{sec:5.4.2} gives for any AI role in a democratic process, and its last condition — a multi-stakeholder review body that includes more than whichever party happens to hold power — is the one that fails first. It fails structurally: an election-monitoring system is procured by an incumbent government, and an incumbent that chooses who audits the count has already answered the question the audit exists to ask.
